Across all experiments, our study investigated how hyperparameters, model size, and reinforcement learning algorithms influence the stability and efficiency of fine-tuning language models for structured query generation. The primary goal was to identify factors that govern stable learning while maximizing performance, focusing on PPO and GRPO algorithms with Qwen2-0.5B and LLaMA 3.2-3B policy models.

The experiments showed that learning rate was the dominant factor affecting stability, with high rates causing extreme oscillations in KL divergence, losses, and entropy collapse. In contrast, reward scaling had minimal impact, while adjustments to KL penalty or PPO epochs influenced exploration and smoothness but did not resolve underlying instability. Larger models, such as LLaMA 3.2-3B, improved baseline rewards and stabilized training dynamics, though learning progress remained limited under the same reward signals.

Switching from PPO to GRPO further enhanced training by smoothing KL divergence, loss curves, and entropy decay. Overall, these findings indicate that learning rate primarily governs stability, model scale improves smoothness and baseline performance, and algorithm choice affects both stability and the efficiency of learning task-specific behaviors.

\subsection{Analysis of Hyperparameters and Their Effects on Stability}
Our experiments clearly demonstrated that hyperparameters play a critical role in shaping the stability and efficiency of PPO-based fine-tuning. Among these, the learning rate emerged as the most influential factor. High learning rates caused extreme fluctuations in KL divergence, value loss, and policy loss, leading to unstable policies and rapid entropy collapse. This instability can be attributed to PPO’s gradient based updates: when the learning rate is too large, even clipped gradients produce abrupt policy shifts, destabilizing the optimization process. Conversely, reducing the learning rate stabilized training by moderating update magnitudes, allowing the policy and value networks to track the reward dynamics more accurately. The trade-off, however, was slower learning progress and early plateauing of average rewards, highlighting the tension between stability and efficiency in RL fine-tuning.

In contrast, reward scaling had minimal effect on training stability. Reducing the reward magnitude by a factor of ten did not prevent KL spikes, loss oscillations, or entropy collapse, indicating that the instability arises primarily from how the optimizer handles gradient updates rather than the absolute reward scale.

Other hyperparameters, such as KL penalty and the number of PPO epochs, influenced the smoothness of optimization and the degree of exploration but were not sufficient to mitigate instability. Lowering the KL coefficient increased policy freedom and slightly improved entropy, while increasing the number of PPO epochs produced smoother loss curves. However, neither adjustment meaningfully accelerated learning, suggesting that these hyperparameters mainly shape optimization dynamics rather than directly controlling learning rate sensitivity.

Overall, these findings emphasize that careful tuning of the learning rate should be the first priority when stabilizing PPO-based fine-tuning, while adjustments to KL penalty or epochs can refine exploration behavior and optimization smoothness without resolving fundamental instability.

\subsection{Effect of Model Architecture and Size}
The choice of model architecture and scale had a noticeable impact on both stability and learning dynamics. Comparing Qwen2-0.5B with LLaMA 3.2-3B, larger models demonstrated smoother KL divergence curves, lower and more stable value losses, and more controlled entropy decay. This indicates that increased model capacity improves representational stability, enabling the policy to track rewards more accurately and avoid abrupt updates.

Despite these advantages, larger models did not guarantee faster or deeper learning under the same reward signal. Even with LLaMA 3.2-3B, the average reward plateaued relatively early, suggesting that learning remained constrained by reward sparsity or insufficient signal strength rather than model capability. This highlights that model scale alone cannot overcome limitations in the reward design or optimization dynamics.

Architectural features may also contribute to these differences. LLaMA’s tokenizer, pretraining corpus, and instruction-tuning regime likely provide a more robust initialization, facilitating smoother policy updates and better alignment with structured query tasks. In contrast, smaller models such as Qwen2-0.5B are more sensitive to hyperparameter choices and gradient noise, making them prone to instability even under identical training conditions. These observations underscore that while larger models improve baseline stability and smoothness.

\subsection{PPO vs GRPO Comparison}
Switching from PPO to GRPO with the LLaMA 3.2-3B policy model and the combined format + recall reward substantially improved both training stability and learning efficiency. GRPO maintained smooth dynamics across all major indicators,  without the extreme oscillations observed in PPO. The entropy declined steadily from around 0.5 to 0.1, reflecting the model’s growing confidence in generating structured Boolean-style queries and rapid learning of the format reward, which quickly reached its maximum limit. KL divergence increased initially during early exploration but stabilized around 1.0–1.2, indicating controlled exploration without the explosive deviations that destabilized PPO.

While GRPO clearly outperformed PPO in overall stability and learning efficiency, the recall component of the reward remained challenging to optimize. Learning progress on recall was slower and struggled to improve consistently, suggesting that GRPO’s group-based averaging preserves stability but can diffuse fine-grained reward signals. Despite this, GRPO still achieved higher overall learning efficiency than PPO, converging faster, avoiding large oscillations, and maintaining balanced policy updates throughout training.

These results highlight that GRPO is particularly effective for tasks requiring stable, reproducible learning of structured outputs, while performance on more detailed or sparse rewards like recall may be limited by reward design rather than algorithmic capability.

\subsection{Implications for Stable RL Fine-Tuning}
The findings of this study provide several important insights into how RL can be applied more reliably for fine-tuning large language models. First, the learning rate emerged as the most critical factor influencing stability. Small adjustments to this parameter dramatically changed training dynamics, from extreme oscillations and value loss explosions at higher rates to smooth, reproducible optimization when reduced. This finding suggests that future fine-tuning efforts should prioritize careful learning rate selection above all other hyperparameter tuning, as stability ultimately depends on controlling the magnitude of policy updates rather than modifying reward scales or other secondary parameters. Practitioners should begin with conservative values and only increase gradually while monitoring KL divergence and value loss trends. Stability should be prioritized over short-term reward gains, as unstable updates can cause irreversible policy collapse.

Second, model size also showed a strong effect on training stability. Larger models offer noticeably smoother optimization and stronger baseline performance, but they also require careful management of compute and memory resources. This indicates that increased capacity helps the model better represent the reward landscape and adapt to reinforcement signals, leading to more stable training. However, larger models did not necessarily learn faster, suggesting that scale improves robustness but not learning efficiency on its own. At last, select algorithms that promote smooth and consistent updates. GRPO demonstrated more stable optimization than PPO, particularly for structured text generation tasks. For practitioners, this suggests prioritizing algorithms that leverage group-based feedback or relative comparisons when stability is a primary concern.

In summary, researchers fine-tuning LLMs with RL should adopt a stability first strategy: begin with a small learning rate, use larger model when resources are available, at least model with 3B parameters, and prefer algorithms like GRPO for reliable optimization. These practices can help ensure smoother, more reproducible training for future RL-based language model applications.

\subsection{Limitations and Future Work}
This study was limited in scope by the selection of models and reward configurations. Only two policy models were tested, Qwen2-0.5B and LLaMA 3.2-3B, each representing a single size from their respective model families. Evaluating additional sizes, such as Qwen2-1.5B, Qwen2-7B, or LLaMA 3.2-1B, would provide a more comprehensive understanding of how model scale influences training stability, learning dynamics, and retrieval performance within the same architecture.

Another limitation lies in the reward design. The original DeepRetrieval framework~\cite{jiang2025deepretrieval} used Recall@3000 to measure retrieval effectiveness, whereas this study was restricted to Recall@50 due to computational and time constraints. A larger retrieval window like Recall@3000 could provide denser and more informative reward signals, allowing the model to learn faster and with clearer feedback on retrieval performance.

Future work should therefore focus on systematically exploring model size effects, extending reward metrics to larger recall cutoffs, and conducting a more detailed analysis of hyperparameters, particularly learning rate sensitivity and its interaction with KL control and optimization frequency. Such investigations would help establish stronger general principles for stable and efficient reinforcement learning fine-tuning of large language models.